\documentclass[10pt,a4paper]{amsart}
\usepackage{amsmath,amssymb,amsthm,graphicx}
\usepackage{verbatim}
\newenvironment{reference}{\comment}{\endcomment}
%\newenvironment{reference}{}{}
\newenvironment{slogan}{\comment}{\endcomment}
\newenvironment{history}{\comment}{\endcomment}

% For commutative diagrams we use Xy-pic
\usepackage[all]{xy}

% We use 2cell for 2-commutative diagrams.
\xyoption{2cell}
% Theorem environments.
%
\theoremstyle{plain}
\newtheorem{theorem}[subsection]{Theorem}
\newtheorem{proposition}[subsection]{Proposition}
\newtheorem{lemma}[subsection]{Lemma}

\theoremstyle{definition}
\newtheorem{definition}[subsection]{Definition}
\newtheorem{example}[subsection]{Example}
\newtheorem{exercise}[subsection]{Exercise}
\newtheorem{situation}[subsection]{Situation}

\theoremstyle{remark}
\newtheorem{remark}[subsection]{Remark}
\newtheorem{remarks}[subsection]{Remarks}

\numberwithin{equation}{subsection}

% Macros
%
\def\lim{\mathop{\mathrm{lim}}\nolimits}
\def\colim{\mathop{\mathrm{colim}}\nolimits}
\def\Spec{\mathop{\mathrm{Spec}}}
\def\Hom{\mathop{\mathrm{Hom}}\nolimits}
\def\Ext{\mathop{\mathrm{Ext}}\nolimits}
\def\SheafHom{\mathop{\mathcal{H}\!\mathit{om}}\nolimits}
\def\SheafExt{\mathop{\mathcal{E}\!\mathit{xt}}\nolimits}
\def\Sch{\mathit{Sch}}
\def\Mor{\mathop{\mathrm{Mor}}\nolimits}
\def\Ob{\mathop{\mathrm{Ob}}\nolimits}
\def\Sh{\mathop{\mathit{Sh}}\nolimits}
\def\NL{\mathop{N\!L}\nolimits}
\def\CH{\mathop{\mathrm{CH}}\nolimits}
\def\proetale{{pro\text{-}\acute{e}tale}}
\def\etale{{\acute{e}tale}}
\def\QCoh{\mathit{QCoh}}
\def\Ker{\mathop{\mathrm{Ker}}}
\def\Im{\mathop{\mathrm{Im}}}
\def\Coker{\mathop{\mathrm{Coker}}}
\def\Coim{\mathop{\mathrm{Coim}}}

% Boxtimes
%
\DeclareMathSymbol{\boxtimes}{\mathbin}{AMSa}{"02}

%
% Macros for moduli stacks/spaces
%
\def\QCohstack{\mathcal{QC}\!\mathit{oh}}
\def\Cohstack{\mathcal{C}\!\mathit{oh}}
\def\Spacesstack{\mathcal{S}\!\mathit{paces}}
\def\Quotfunctor{\mathrm{Quot}}
\def\Hilbfunctor{\mathrm{Hilb}}
\def\Curvesstack{\mathcal{C}\!\mathit{urves}}
\def\Polarizedstack{\mathcal{P}\!\mathit{olarized}}
\def\Complexesstack{\mathcal{C}\!\mathit{omplexes}}
% \Pic is the operator that assigns to X its picard group, usage \Pic(X)
% \Picardstack_{X/B} denotes the Picard stack of X over B
% \Picardfunctor_{X/B} denotes the Picard functor of X over B
\def\Pic{\mathop{\mathrm{Pic}}\nolimits}
\def\Picardstack{\mathcal{P}\!\mathit{ic}}
\def\Picardfunctor{\mathrm{Pic}}
\def\Deformationcategory{\mathcal{D}\!\mathit{ef}}
\usepackage{cleveref}
\crefname{originalreference}{original reference}{original references}
\Crefname{originalreference}{Original reference}{Original references}
\makeatletter
\newlabel{definition-big-proetale-site}{{098G}{source}{}{sourcepointer-1}{}}
\newlabel{definition-big-proetale-site@cref}{{[originalreference][0][]098G}{[1][0][]source}{}{}{}}
\makeatother
\crefname{originalreference}{original reference}{original references}
\Crefname{originalreference}{Original reference}{Original references}
\makeatletter
\newlabel{lemma-w-contractible-is-weakly-contractible}{{098H}{source}{}{sourcepointer-2}{}}
\newlabel{lemma-w-contractible-is-weakly-contractible@cref}{{[originalreference][0][]098H}{[1][0][]source}{}{}{}}
\makeatother
\crefname{originalreference}{original reference}{original references}
\Crefname{originalreference}{Original reference}{Original references}
\makeatletter
\newlabel{proposition-find-w-contractible}{{0983}{source}{}{sourcepointer-3}{}}
\newlabel{proposition-find-w-contractible@cref}{{[originalreference][0][]0983}{[1][0][]source}{}{}{}}
\makeatother
\crefname{originalreference}{original reference}{original references}
\Crefname{originalreference}{Original reference}{Original references}
\makeatletter
\newlabel{remark-size-w-contractible}{{0984}{source}{}{sourcepointer-4}{}}
\newlabel{remark-size-w-contractible@cref}{{[originalreference][0][]0984}{[1][0][]source}{}{}{}}
\makeatother
\crefname{originallemma}{original lemma}{original lemmas}
\Crefname{originallemma}{Original lemma}{Original lemmas}
\makeatletter
\newlabel{sets-lemma-bound-affine}{{000P}{source}{}{sourcepointer-5}{}}
\newlabel{sets-lemma-bound-affine@cref}{{[originallemma][0][]000P}{[1][0][]source}{}{}{}}
\makeatother
\crefname{originallemma}{original lemma}{original lemmas}
\Crefname{originallemma}{Original lemma}{Original lemmas}
\makeatletter
\newlabel{sets-lemma-bound-finite-type}{{04T7}{source}{}{sourcepointer-6}{}}
\newlabel{sets-lemma-bound-finite-type@cref}{{[originallemma][0][]04T7}{[1][0][]source}{}{}{}}
\makeatother
\crefname{originallemma}{original lemma}{original lemmas}
\Crefname{originallemma}{Original lemma}{Original lemmas}
\makeatletter
\newlabel{sets-lemma-construct-category}{{000J}{source}{}{sourcepointer-7}{}}
\newlabel{sets-lemma-construct-category@cref}{{[originallemma][0][]000J}{[1][0][]source}{}{}{}}
\makeatother
\crefname{originallemma}{original lemma}{original lemmas}
\Crefname{originallemma}{Original lemma}{Original lemmas}
\makeatletter
\newlabel{sets-lemma-coverings-site}{{000X}{source}{}{sourcepointer-8}{}}
\newlabel{sets-lemma-coverings-site@cref}{{[originallemma][0][]000X}{[1][0][]source}{}{}{}}
\makeatother
\title{Every object of the big pro-etale site is covered by spectra of w-contractible rings}
\author{The Stacks Project Authors}
\date{stacks-project:proetale.tex:local-snapshot-20261012; GNU FDL 1.2 or later}
\begin{document}
\maketitle
\noindent\textit{Editorial note.} All delivered mathematical text is the author's own, copied from the pinned Stacks Project chapter file proetale.tex; no statement or proof step was rewritten, reordered or removed.\par\medskip
\noindent\textit{Editorial note.} The counted proof is the author's complete argument: the construction of the covering through faithfully flat ind-etale extensions by w-contractible rings, the explicit cardinality estimate showing that the pieces are objects of the big pro-etale site, and the conclusion of weak contractibility.\par\medskip
\noindent\textit{Editorial note.} The label lemma-get-many-weakly-contractible names the counted statement, and the proof environment immediately following it in the source is delivered verbatim as its proof.\par\medskip
\noindent\textit{Editorial note.} The counted statement is carried by the author's own lemma environment in the chapter source.\par\medskip
\noindent\textit{Editorial note.} Pre-open gate: python operations/agent-authoring/tools/precard\_gate.py --source proetale.tex --label lemma-get-many-weakly-contractible --start 3086 --end 3127 --local-id 3586 returned RESULT=OK (L3586 free, occupancy FREE, public interval clear). Candidate taken from sources/stacks-proetale/candidates-halfA.json.\par\medskip
\noindent\textit{Editorial note.} The proof opens by addressing set-theoretic concerns explicitly ('For those readers who do not care about set-theoretical issues this lemma is a trivial consequence ... Here are the details.'); the details that follow are the argument, so the statement is counted.\par\medskip
\noindent\textit{Editorial note.} The statement and proof are copied from the pinned chapter file; the cited inputs are kept as external pointers with their canonical Stacks tags.\par\medskip
\noindent\textit{Allowed standard prerequisites (editorial).}\begin{itemize}
\item Existence of faithfully flat ind-etale extensions by w-contractible rings. Same chapter, the proposition producing a faithfully flat ind-etale ring map from a ring to a w-contractible ring. Supplies the maps A\_i -> D\_i with D\_i w-contractible used to build the covering.
\item Equivalence between w-contractibility and weak contractibility. Same chapter, the lemma identifying w-contractible affine schemes with weakly contractible objects of the big pro-etale site. Supplies the conclusion that the spectra of the w-contractible rings are weakly contractible.
\item The construction of the big pro-etale site and its size function. Same chapter, the definition of the big pro-etale site including the size bound, and the lemma describing which schemes enter the constructed category. Supply the criterion by which Spec(D\_i) is isomorphic to an object of the big pro-etale site, namely the bound on the cardinality of D\_i.
\item Estimates for the size of affine schemes and of w-contractible rings. Sets, the lemmas bounding the size of affine schemes and of schemes of finite type, and Same chapter, the remark on the size of a w-contractible ring. Supply the estimate of the cardinality of D\_i in terms of the size of T, which is within the chosen bound.
\end{itemize}\medskip
\ifdefined\hypertarget\hypertarget{sourcepointer-1}{}\fi
\noindent\textit{Source reference.} Original reference 098G: sources/stacks-proetale/pinned-original-source.tex, source lines 2465--2465. Prerequisite pointer invoked at proof line(s) 3089; it names a standard fact of the same chapter that the counted argument uses, and every step of the counted proof built on it is retained verbatim in the delivered body.\par\medskip
\ifdefined\hypertarget\hypertarget{sourcepointer-2}{}\fi
\noindent\textit{Source reference.} Original reference 098H: sources/stacks-proetale/pinned-original-source.tex, source lines 3027--3027. Prerequisite pointer invoked at proof line(s) 3099; it names a standard fact of the same chapter that the counted argument uses, and every step of the counted proof built on it is retained verbatim in the delivered body.\par\medskip
\ifdefined\hypertarget\hypertarget{sourcepointer-3}{}\fi
\noindent\textit{Source reference.} Original reference 0983: sources/stacks-proetale/pinned-original-source.tex, source lines 2135--2135. Prerequisite pointer invoked at proof line(s) 3100; it names a standard fact of the same chapter that the counted argument uses, and every step of the counted proof built on it is retained verbatim in the delivered body.\par\medskip
\ifdefined\hypertarget\hypertarget{sourcepointer-4}{}\fi
\noindent\textit{Source reference.} Original reference 0984: sources/stacks-proetale/pinned-original-source.tex, source lines 2171--2171. Prerequisite pointer invoked at proof line(s) 3124; it names a standard fact of the same chapter that the counted argument uses, and every step of the counted proof built on it is retained verbatim in the delivered body.\par\medskip
\ifdefined\hypertarget\hypertarget{sourcepointer-5}{}\fi
\noindent\textit{Source reference.} Original lemma 000P: sources/stacks-chow/tags.txt, source lines 34--34. Prerequisite pointer invoked at proof line(s) 3121; it names a standard fact of a sibling chapter that the counted argument uses, and every step of the counted proof built on it is retained verbatim in the delivered body.\par\medskip
\ifdefined\hypertarget\hypertarget{sourcepointer-6}{}\fi
\noindent\textit{Source reference.} Original lemma 04T7: sources/stacks-chow/tags.txt, source lines 6245--6245. Prerequisite pointer invoked at proof line(s) 3122; it names a standard fact of a sibling chapter that the counted argument uses, and every step of the counted proof built on it is retained verbatim in the delivered body.\par\medskip
\ifdefined\hypertarget\hypertarget{sourcepointer-7}{}\fi
\noindent\textit{Source reference.} Original lemma 000J: sources/stacks-chow/tags.txt, source lines 28--28. Prerequisite pointer invoked at proof line(s) 3119; it names a standard fact of a sibling chapter that the counted argument uses, and every step of the counted proof built on it is retained verbatim in the delivered body.\par\medskip
\ifdefined\hypertarget\hypertarget{sourcepointer-8}{}\fi
\noindent\textit{Source reference.} Original lemma 000X: sources/stacks-chow/tags.txt, source lines 42--42. Prerequisite pointer invoked at proof line(s) 3113; it names a standard fact of a sibling chapter that the counted argument uses, and every step of the counted proof built on it is retained verbatim in the delivered body.\par\medskip
\begin{lemma}
\label{lemma-get-many-weakly-contractible}
Let $\Sch_\proetale$ be a big pro-\'etale site as in
Definition \ref{definition-big-proetale-site}.
For every object $T$ of $\Sch_\proetale$ there exists
a covering $\{T_i \to T\}$ in $\Sch_\proetale$
with each $T_i$ affine and the spectrum of a w-contractible
ring. In particular, $T_i$ is weakly contractible in $\Sch_\proetale$.
\end{lemma}

\begin{proof}
For those readers who do not care about set-theoretical issues
this lemma is a trivial consequence of
Lemma \ref{lemma-w-contractible-is-weakly-contractible} and
Proposition \ref{proposition-find-w-contractible}.
Here are the details.
Choose an affine open covering $T = \bigcup U_i$. Write $U_i = \Spec(A_i)$.
Choose faithfully flat, ind-\'etale ring maps $A_i \to D_i$
such that $D_i$ is w-contractible as in
Proposition \ref{proposition-find-w-contractible}.
The family of morphisms $\{\Spec(D_i) \to T\}$ is a
pro-\'etale covering.
If we can show that $\Spec(D_i)$ is isomorphic to an object, say $T_i$,
of $\Sch_\proetale$, then $\{T_i \to T\}$ will be combinatorially
equivalent to a covering of $\Sch_\proetale$ by the construction
of $\Sch_\proetale$ in Definition \ref{definition-big-proetale-site}
and more precisely the application of
Sets, Lemma \ref{sets-lemma-coverings-site} in the last step.
To prove $\Spec(D_i)$ is isomorphic to an object of
$\Sch_\proetale$, it suffices to prove that
$|D_i| \leq Bound(\text{size}(T))$ by the construction
of $\Sch_\proetale$ in Definition \ref{definition-big-proetale-site}
and more precisely the application of
Sets, Lemma \ref{sets-lemma-construct-category} in step (3).
Since $|A_i| \leq \text{size}(U_i) \leq \text{size}(T)$
by Sets, Lemmas \ref{sets-lemma-bound-affine} and
\ref{sets-lemma-bound-finite-type} we get
$|D_i| \leq \kappa^{2^{2^{2^\kappa}}}$ where $\kappa = \text{size}(T)$
by Remark \ref{remark-size-w-contractible}.
Thus by our choice of the function $Bound$ in
Definition \ref{definition-big-proetale-site} we win.
\end{proof}

\end{document}
